\chapter{Fazit und Ausblick}
\label{ch:chapter07}

In diesem Kapitel werden die wesentlichen Erkenntnisse der Arbeit zusammengefasst, der wissenschaftliche Beitrag eingeordnet, die 
Limitationen diskutiert und ein Ausblick auf mögliche weiterführende Fragestellungen gegeben.

\section{Zusammenfassung der Arbeit}

Die vorliegende Arbeit analysierte die sprachspezifischen Performanceveränderungen durch OpenTelemetry in AWS-Lambda-Funktionen für Java, 
Python und Node.js. Dazu wurden fünf praxisnahe Microbenchmarks jeweils in instrumentierter und nicht instrumentierter Variante 
unter kontrollierten Bedingungen gemessen. Die Auswertung erfolgte auf Basis von Bootstrap-Resampling und Konfidenzintervallen.

Die Ergebnisse zeigen konsistent, dass Python den geringsten relativen Overhead hat, gefolgt von Java, während Node.js vor allem 
wegen hoher Init-Durations bei Kaltstarts am schwächsten abschneidet. Auch in der absoluten Performance mit OpenTelemetry liegt 
Python vorn, und die Kostenanalyse bestätigt Python in vier von fünf Benchmarks als günstigste Option. Die beobachteten negativen 
Overheads bei HTTP-Kaltstarts in Node.js und Python sind statistisch signifikant, praktisch jedoch vernachlässigbar.

\section{Wissenschaftlicher Beitrag}

Die Arbeit leistet einen empirischen Beitrag zur Serverless-Forschung, indem sie erstmals den sprachspezifischen Overhead von 
OpenTelemetry auf AWS Lambda systematisch quantifiziert. In Kontrast zu früheren Studien, die entweder nur eine Sprache 
(Eder et al.~\cite{10254754}) oder nur die absolute Performance ohne Tracing (Rodriguez et al.~\cite{Rodriguez2024Evaluation}) 
untersuchten, kombiniert diese Arbeit beide Perspektiven. Die Verwendung von fünf realitätsnahen Microbenchmarks, die strikte 
Kontrolle der Konfigurationsparameter und die statistische Absicherung durch Hypothesentests erhöhen die Aussagekraft der Ergebnisse. 
Zudem werden alle Skripte und Rohdaten in einem öffentlichen GitHub-Repository bereitgestellt, was die Reproduzierbarkeit und Nachnutzung 
durch andere Forschende ermöglicht.

\section{Limitationen der Arbeit}

Die vorliegende Studie ist durch mehrere Faktoren eingeschränkt. Alle Experimente wurden ausschließlich auf AWS Lambda durchgeführt, 
sodass sich die Erkenntnisse nicht ohne Weiteres auf andere FaaS-Plattformen wie Azure Functions oder Google Cloud Functions übertragen 
lassen. Zudem wurde nur eine Speichergröße von 1024 MB verwendet, obwohl der verfügbare Speicher Laufzeit und Kaltstartverhalten 
erheblich beeinflussen kann, insbesondere bei Java~\cite{Rodriguez2024Evaluation}. Auch die Analyse bleibt auf Head Sampling mit einer 
Rate von 100\,\% beschränkt. Schließlich decken die fünf Microbenchmarks zwar zentrale Anwendungsfälle ab, bilden aber weder 
rechenintensive noch stark I/O-lastige Szenarien vollständig ab. Insgesamt verringern diese Einschränkungen nicht den Wert der 
Ergebnisse, zeigen jedoch, dass weitere Untersuchungen für andere Plattformen, Konfigurationen und Workloads sinnvoll wären.

\section{Ausblick auf zukünftige Arbeiten}

Aus den Limitationen dieser Arbeit ergeben sich mehrere Forschungsrichtungen. Empfehlenswert sind zunächst vergleichende Untersuchungen 
auf anderen FaaS-Plattformen wie Azure Functions oder Google Cloud Functions, um die Allgemeinheit der beobachteten sprachspezifischen 
Overheads zu prüfen. Ebenso sollte analysiert werden, wie sich der sprachspezifische Overhead bei anderen Plattformen verhält, etwa durch 
Vergleiche mit leichteren Frameworks wie Zipkin oder nativen Alternativen wie AWS X-Ray. Darüber hinaus ist offen, wie sich 
unterschiedliche Sampling-Strategien, insbesondere Tail Sampling oder niedrigere Sampling-Raten, auf Performance und Kosten auswirken. 
Für Node.js erscheint zudem eine vertiefte Ursachenanalyse der hohen Init-Duration-Overheads sinnvoll. Das gleiche gilt für die 
auffälligen negativen Overheads bei HTTP-Kaltstarts, die in kontrollierten Experimenten weiter untersucht werden sollten. 
Schließlich könnte eine Ausweitung auf weitere Programmiersprachen wie Go, .NET (C\#) oder Rust die Generalisierbarkeit der Ergebnisse erhöhen. 
\section{Abschließende Bemerkung}

Abschließend zeigt die Arbeit, dass die Programmiersprache einen deutlichen Einfluss auf den durch OpenTelemetry verursachten Performance-Overhead in serverlosen Umgebungen hat. Python ist für die meisten synchronen Anwendungsfälle die beste Wahl, während Java in einzelnen asynchronen Szenarien konkurrenzfähig sein kann. Node.js ist wegen der hohen Kaltstart-Overheads nur eingeschränkt empfehlenswert. Damit bietet die Arbeit eine Grundlage für die Sprachauswahl bei observierbaren serverlosen Systemen.